% Catalog architecture — describe the three layers (explicit / derived /
% cited), the scheme JSON format, and the lineage records. Should
% be technical enough for someone to be able to recreate the catalog
% from scratch.

The catalog's source of truth is a tree of \emph{per-scheme} JSON files
under \code{src/main/resources/schemes/section\(N\)/}, bucketed by the
maximum of the three shape dimensions. Each file is one algorithm ---
one shape, one field, one commutativity tag. We do not store every
scheme the search can build, only the ones worth keeping:
\begin{itemize}[itemsep=2pt]
  \item the \emph{elected} best derived scheme for each shape --- we aim
    for at least one per shape, the current best-in-scope construction;
  \item \emph{atoms} worth cataloguing in their own right --- a low
    rank, a low addition count, historical provenance, or any other
    reason to preserve the specimen;
  \item schemes kept for their \emph{ingredient} value: a serendipity or
    projection scheme may not itself hold the best rank at its shape yet
    be actively consumed as a sub-product by another shape's best
    scheme, so we retain it.
\end{itemize}
On top of this tree sit a few \emph{aggregated} JSON files under
\code{docs/}, \emph{generated} from it: the manifest
\code{docs/catalog.json} (one row per elected scheme, consumed by the
web browser and the comparison tooling), alongside
\code{docs/cited-bounds.json} and
\code{docs/derived-from-cited-bounds.json}.

\subsection{Three layers}

\paragraph{Explicit schemes.} JSON files that carry the full \((U,
V, W)\) factor matrices; every explicit scheme passes a randomised
spot-check at load time. The filename is a pure cosmetic label
(shape, rank, source or derivation note, short content hash): every
property --- field, addition count, commutativity --- is read from
the JSON content, never parsed from the name. The hash, the first
seven hex digits of a content hash of the factor matrices, gives
each scheme a stable identity independent of how it was named
upstream.

\paragraph{Derived bounds.} A separate JSON file
\code{docs/derived-from-cited-bounds.json} records rank claims that our own
constructors can re-derive on demand from formulas (Waksman 1970,
Rosowski 2019 Theorem 2/3, Pan 1980 trilinear aggregation, Hopcroft--Kerr
1971 \(\nmpshape{2}{b}{c}\)). Generated by
\code{GenerateDerivedBounds.java}; rebuilds whenever the underlying
formula changes. This is the layer that grows as we ``move'' bounds
from the cited tier to the derived tier by reproducing the original
construction.

\paragraph{Cited bounds.} The failover tier. A central JSON file
\code{docs/cited-bounds.json} records rank claims whose explicit
factor matrices we do not (yet) hold and which no constructor of ours
re-derives --- so we neither materialise nor re-derive them, but still
trust and attribute the published rank. Each entry has shape, field,
commutativity, rank, addition count (or \code{null}), source
citation, and discovery status. This lets us show the result in
comparison tables without pretending to have the scheme. The cited
layer is generated by \code{GenerateCitedBounds.java}.

\subsection{Lineage records}

The \emph{lineage} of a scheme is the exact sequence of composition
steps that produced it, and it is \textbf{replayable}: a
\code{LineageReplayer} reconstructs the full concrete \((U, V, W)\)
factor matrices from the lineage alone. This is essential --- it lets
the catalog store a large scheme as a lineage-only \emph{stub} (no
factor matrices on disk) and rebuild them on demand, and it makes every
composed scheme reproducible from its building blocks rather than
trusted as opaque numbers. The lineage is a DAG whose leaf nodes are
atoms and whose internal nodes are the derivation operators of
Section~\ref{sec:strategies}, one node type per operator. The lineage
is serialised twice: as a
human-readable function-call form (e.g.\
\code{ConcatRight(KronProduct(strassen, winograd), SAME0)})\, and as
a machine-readable JSON DAG with shared-subtree references.

The lineage is load-bearing for three reasons:
\begin{enumerate}[itemsep=2pt]
  \item \textbf{Independent audit.} Because replay is independent of
    any stored matrices, re-deriving from leaves and comparing catches
    silent corruption of an on-disk \((U, V, W)\).
  \item \textbf{Search-without-materialisation.} The exhaustive derivation
    search (Section \ref{sec:search}) writes a strategy descriptor
    -- effectively a lineage stub -- before deciding whether to
    materialise. This lets the search iterate without paying for
    materialisation.
  \item \textbf{Attribution archaeology.} When a bulk-imported
    catalog re-encodes a scheme originally due to an earlier author,
    the lineage records the importing source while a separate
    attribution field (\code{attribution\_for\_rank}) records the
    earliest source that established the bound. Comparison tables
    pull from the latter.
\end{enumerate}

\subsection{Verification}

Four verification modes run independently of the search loop.
\begin{description}[itemsep=2pt]
  \item[Spot-check.] \code{Verifier.passesRandomMatmulSpotCheck}
    samples random \(A\), \(B\), evaluates the scheme via its
    \((U, V, W)\), and compares against the direct product ---
    \(O(S\,r\,(nm{+}mp{+}np))\) for \(S\) samples. Cheap enough to run
    at scheme load time, after every materialisation, and as a nightly
    catalog-wide sweep; probabilistic, but a few samples over a field
    make a false pass vanishingly unlikely.
  \item[Symbolic.] \code{SymbolicVerifier} checks the trilinear
    identity \emph{exactly} (an algebraic equality, no floating-point
    epsilon) and additionally rejects any coefficient outside the
    declared field. Conclusive, but \(O(r\,nm\,mp\,np)\) --- the full
    tensor, \(O(r\,n^6)\) in the cubic case --- so it is reserved for
    the small-dimension schemes where that is affordable.
  \item[Lineage replay.] \code{LineageReplayer} re-constructs the
    scheme from its lineage record, requiring the same
    \((U, V, W)\). Catches divergence between the on-disk scheme and
    the formula it claims to instantiate.
  \item[Field widening.] \code{FieldWideningSweep} re-tags a scheme
    claimed for a wide field down to the strictest field its
    coefficients inhabit (e.g.\ an \(\Reals\)-tagged scheme with all
    coefficients in \(\Rationals\) is really \(\Rationals\)). Within
    characteristic~0 the coefficients \emph{guarantee} the narrower
    tag --- matmul-validity is a system of integer polynomial
    identities, unchanged by restricting to a subfield --- so this is a
    certain re-tag, not a heuristic.
\end{description}
